\documentclass[12pt]{article}

\usepackage{exsheets}

%\usepackage[pdf]{pstricks}
%\usepackage{pstricks-add}

%\usepackage{tikz} 
\usepackage[fleqn]{amsmath}
\usepackage{amssymb}
\usepackage{amstext}  
\usepackage{amsthm} 
\usepackage{array}
\usepackage{mathtools}
\mathtoolsset{showonlyrefs} 
\usepackage{graphicx,hyperref}
\usepackage{listings}
\usepackage{indentfirst}
\usepackage[top=1.5cm, bottom=1.5cm, left=.5in, right=.5in]{geometry}
\usepackage{enumerate}
\usepackage{multicol}

\newcommand{\ssol}{\vspace{3em}}
\newcommand{\lsol}{\vspace{10em}}
\newcommand{\blnk}{{\underline {\hspace{1.5in}}}}

%\usepackage{mcode}

\setcounter{secnumdepth}{0} % Turns off automatic numbering of sections.

\SetupExSheets{headings=block}



\begin{document}
Justin Boyer\\
Master's Thesis Project Proposal

\section{Overview}

I will create a document targeted at teachers of secondary mathematics that explains the similarities and differences between the polynomials with integer coeffecients and the integers.  I will leverage the explanation of these similarites and differences to provide the reader with a way to teach arithmetic operations on the polynomials with integer coeffecients by connecting them to arithmetic on the integers.  I will then expand the scope of the project to include polynomials with rational coeffecients, and vectors with integer elements.  I will then compare and contrast clock arithmetic in each ring.  I will close the paper with an application in coding theory in which High School students could put their new knowledge to the test.  The students will use to different representations to multiply n-nomials (where $n \in \mathbb{N}$) and then apply coding theory (an application of abstract algebra) to their practices to correct their errors.

\section{Outline}




\end{document}


\begin{pspicture}
\psgrid[unit=.05\linewidth, subgriddiv=1, linewidth=.01 gridlabels=0pt, gridcolor=gray, griddots=10](0,0)(-10,-10)(10,10) 
\psaxes[unit=.05\linewidth, linewidth=.1 labelFontSize=\scriptscriptstyle]{<->}(0,0)(-10.5,-10.5)(10.5,10.5)
\end{pspicture}


\begin{tikzpicture}[scale=.5]	
\draw[help lines] (0,0) grid (20,20);
\end{tikzpicture}


